We consider an increasing sequence $A=\{a_1<a_2<\cdots\}\subseteq\mathbb{N}$ and its subset-sum set
\[
P(A)=\left\{\sum_{n\in B} n:\ B\subseteq A\ \text{finite}\right\}.
\]
A sequence $A$ is \emph{complete} if every sufficiently large integer lies in $P(A)$.
The problem asks whether there exists a sequence $A$ such that:
\begin{enumerate}
\item[(1)] For every finite subset $B\subseteq A$, the set $A\setminus B$ is complete.
\item[(2)] For every infinite subset $B\subseteq A$, the set $A\setminus B$ is \emph{not} complete.
\end{enumerate}
The statement reports that Graham exhibited an example with these properties.
It further asks: if additionally $a_{n+1}/a_n\ge 1+\varepsilon$ for all $n$, must one have
\[
\lim_{n\to\infty}\frac{a_{n+1}}{a_n}=\frac{1+\sqrt5}{2}=:\varphi\ ?
\]
It also notes that if $a_{n+1}/a_n>\varphi$ for all $n$, then property (2) is ``easy''.
